% Appendix: Read Directions — Definition Levers
% Standalone mathematical reference for the design space of read-direction definitions.
\PassOptionsToPackage{unicode}{hyperref}
\PassOptionsToPackage{hyphens}{url}
\documentclass[
]{article}
\usepackage{xcolor}
\usepackage{amsmath,amssymb}
\setcounter{secnumdepth}{-\maxdimen} % remove section numbering
\usepackage{iftex}
\ifPDFTeX
  \usepackage[T1]{fontenc}
  \usepackage[utf8]{inputenc}
  \usepackage{textcomp} % provide euro and other symbols
\else % if luatex or xetex
  \usepackage{unicode-math} % this also loads fontspec
  \defaultfontfeatures{Scale=MatchLowercase}
  \defaultfontfeatures[\rmfamily]{Ligatures=TeX,Scale=1}
\fi
\usepackage{lmodern}
% Use upquote if available, for straight quotes in verbatim environments
\IfFileExists{upquote.sty}{\usepackage{upquote}}{}
\IfFileExists{microtype.sty}{% use microtype if available
  \usepackage[]{microtype}
  \UseMicrotypeSet[protrusion]{basicmath} % disable protrusion for tt fonts
}{}
\makeatletter
\@ifundefined{KOMAClassName}{% if non-KOMA class
  \IfFileExists{parskip.sty}{%
    \usepackage{parskip}
  }{% else
    \setlength{\parindent}{0pt}
    \setlength{\parskip}{6pt plus 2pt minus 1pt}}
}{% if KOMA class
  \KOMAoptions{parskip=half}}
\makeatother
\setlength{\emergencystretch}{3em} % prevent overfull lines
\providecommand{\tightlist}{%
  \setlength{\itemsep}{0pt}\setlength{\parskip}{0pt}}
\usepackage{bookmark}
\IfFileExists{xurl.sty}{\usepackage{xurl}}{} % add URL line breaks if available
\urlstyle{same}
\hypersetup{
  hidelinks}

\newcommand{\dmodel}{d_{\mathrm{model}}}
\newcommand{\dhead}{d_{\mathrm{head}}}
\DeclareMathOperator*{\argmax}{arg\,max}

\author{}
\date{}

\begin{document}

\section{Appendix: Read Directions — Definition Levers}\label{read-direction-levers}

A \emph{read direction} is an input-side direction of the residual stream that the OV
circuits of the selected function-vector heads map onto the task imitation space: the
direction the circuit ``reads'' in order to write the task content it outputs. Given a
head $h$ with OV circuit $W^h_O W^h_V \in \mathbb{R}^{\dmodel \times \dmodel}$, its mean
task activation $h_A \in \mathbb{R}^{\dmodel}$, the selected head set $H$, and the
function vector $v_A = \sum_{h \in H} h_A$, a read direction is some vector
$r_A \in \mathbb{R}^{\dmodel}$ whose image under the circuit is aligned with the task
content. For GPT-J-6B, $\dmodel = 4096$ and $\dhead = 256$, so each single-head circuit
$W^h_O W^h_V$ has rank at most $256$ and a kernel of dimension at least $3840$.

There is more than one defensible way to make ``aligned with the task content'' precise,
and the choices interact. Rather than committing to a single definition, we state the
design space as four levers and sweep over the resulting definitions, selecting among
them empirically. Throughout, ``target'' denotes the vector the read direction should
reproduce through the circuit --- $h_A$ for a single head, $v_A$ (or its per-prompt
analogue $v^j_A$) for the aggregated constructions. The choice of task-level versus
per-prompt target is orthogonal to the levers below and is held fixed within any sweep.

\subsection{Lever 1 — Optimisation metric: dot product vs.\ cosine
similarity}\label{lever-1-metric}

The first choice is what ``the input direction whose image is maximally aligned with the
target'' means. Let $W$ stand for the relevant circuit (a single head's $W^h_O W^h_V$,
or the summed circuit $M$ of Lever 3) and $t$ for the target.

\begin{enumerate}
\item
  \textbf{Dot product.} Maximise the inner product of the image with the target,
  \begin{equation}
  r \;=\; \argmax_{\|z\| = 1} \; \langle W z,\, t \rangle
  \;\;\propto\;\; W^{\top} t. \tag{1}
  \end{equation}
  The closed form follows from
  $\langle Wz, t\rangle = \langle z, W^{\top} t\rangle$. No kernel constraint is
  needed: $W^{\top} t$ lies in the row space of $W$ automatically. This objective
  rewards output \emph{magnitude}: input directions the circuit transmits strongly are
  favoured even when their image is less purely aligned with $t$.
\item
  \textbf{Cosine similarity.} Maximise the angle-based alignment of the image with the
  target,
  \begin{equation}
  r \;=\; \argmax_{\substack{\|z\| = 1 \\ z \,\perp\, \ker(W)}}
  \; \cos\!\left(W z,\, t\right)
  \;\;\propto\;\; W^{+} t, \tag{2}
  \end{equation}
  where $W^{+}$ is the Moore--Penrose pseudo-inverse and the kernel constraint makes
  the maximiser unique (adding any $u \in \ker(W)$ to a maximiser leaves the output
  unchanged). This objective is indifferent to output magnitude and rewards purity of
  alignment; the price is that $W^{+}$ scales components by $1/\sigma_i$, so
  near-zero singular values can dominate the solution. Lever 2 exists to control this.
\end{enumerate}

The two objectives are genuinely different optimisation problems with different
solutions: the unit-norm constraint fixes the norm of the \emph{input} $z$, but the
cosine additionally normalises by the norm of the \emph{output} $Wz$, which still varies
with $z$. The dot product rewards output magnitude along $t$; the cosine does not.

\subsection{Lever 2 — Inversion: literal pseudo-inverse vs.\
\texorpdfstring{$\tau$}{tau}-truncation (cosine metric only)}\label{lever-2-inversion}

This lever applies when Lever 1 selects the cosine metric; the dot-product solution
involves no inversion, so the lever collapses there.

\begin{enumerate}
\item
  \textbf{Literal pseudo-inverse $W^{+}$.} Invert every nonzero singular value. This is
  the exact maximiser of the cosine objective, but it amplifies whatever the circuit
  barely transmits: components of $t$ along output directions with small $\sigma_i$
  receive weight $1/\sigma_i$ in the solution, so the read direction can be dominated
  by directions of near-zero causal throughput.
\item
  \textbf{$\tau$-truncated pseudo-inverse $W^{+}_{\tau}$.} Invert only the singular
  values $\sigma_i \ge \tau$; equivalently, additionally constrain $z$ to be orthogonal
  to the right-singular subspace with $\sigma_i < \tau$. This trades a small loss in
  achievable cosine for robustness to directions the circuit does not meaningfully
  transmit. The threshold $\tau$ (equivalently, a rank cut $k$) is a hyperparameter of
  the sweep; $\tau \to 0$ recovers the literal pseudo-inverse, and truncation to the
  top singular direction is the opposite extreme.
\end{enumerate}

(A truncated variant of the dot-product solution --- projecting $W^{\top} t$ onto the
top-$k$ right-singular subspace --- is definable but is not part of the present
crossing.)

\subsection{Lever 3 — Aggregation across heads: summed circuit vs.\
per-head-then-sum}\label{lever-3-aggregation}

The function vector is carried by multiple heads, so a single task-level read direction
requires a choice of how the heads combine.

\begin{enumerate}
\item
  \textbf{Summed circuit.} Form the summed OV circuit of the selected heads,
  \begin{equation}
  M \;=\; \sum_{h \in H} W^h_O W^h_V \;\in\; \mathbb{R}^{\dmodel \times \dmodel},
  \tag{3}
  \end{equation}
  take the target $t = v_A$ (or $v^j_A$), and solve the Lever 1 problem once against
  $M$. This treats every head as reading the \emph{same} residual-stream vector $z$,
  even though the heads sit at different layers --- a physical inconsistency accepted
  for tractability. Note that $\operatorname{rank}(M)$ can reach $\dhead\,|H|$, so with
  $|H| \gtrsim \dmodel/\dhead = 16$ heads $M$ can be full rank; the kernel constraint
  of eq.~(2) is then vacuous and $\tau$-truncation (Lever 2) becomes the only
  regulariser.
\item
  \textbf{Per-head, then sum and normalise.} For each head $h \in H$, solve the Lever 1
  problem for that head's own circuit $W^h_O W^h_V$ against that head's own mean
  activation $h_A$, giving a per-head read direction $r^h_A$; then aggregate:
  \begin{equation}
  r_A \;=\; \frac{\sum_{h \in H} r^h_A}{\bigl\|\sum_{h \in H} r^h_A\bigr\|}. \tag{4}
  \end{equation}
  This respects each head's own input geometry but loses cross-head interactions: no
  head can compensate for another's blind spot, unlike in the summed circuit.

  \emph{Sub-choice (3$'$).} Each $r^h_A$ is unit-norm by construction, so the plain sum
  in eq.~(4) weights heads equally. An alternative weights each head by the magnitude
  of its unnormalised solution --- $\|(W^h_O W^h_V)^{+} h_A\|$ under the cosine metric,
  or $\|(W^h_O W^h_V)^{\top} h_A\|$ under the dot product --- before summing.
\end{enumerate}

\subsection{Lever 4 — Final normalisation: unit norm or natural
magnitude}\label{lever-4-normalisation}

The last choice is whether the final $r_A$ is unit-normalised.

\begin{enumerate}
\item
  \textbf{Unit norm.} $r_A$ is a pure direction; any intervention strength is delegated
  entirely to the steering-coefficient sweep over $\alpha$.
\item
  \textbf{Natural magnitude.} Keep the norm the construction produces ---
  $\|W^{+} t\|$ under the cosine metric, $\|W^{\top} t\|$ under the dot product, or the
  norm of the per-head sum in Lever 3.2. These magnitudes carry meaning (for instance,
  $\|M^{+} v_A\|$ is the input scale required to produce a $v_A$-sized output through
  $M$), and they differ across tasks and heads, so this option changes the relative
  dosing across tasks in any comparison at fixed $\alpha$.
\end{enumerate}

This lever is immaterial for purely correlational uses of read directions (the cosine of
$r_A$ against residual activations is scale-invariant); it matters for causal uses ---
injection, ablation, or projection-removal at a fixed intervention strength.

\subsection{The resulting crossing}\label{the-resulting-crossing}

Under the dot-product metric Lever 2 collapses, giving
$1 \times 2 \times 2 = 4$ definitions; under the cosine metric the crossing is
$\{\text{literal},\ \tau\text{-truncated over a grid of } \tau\} \times 2 \times 2$.
Two points of the crossing correspond to the definitions used in the main text: the
single-head read direction $r^h_A$ (cosine, literal pseudo-inverse, per-head, unit norm)
and the summed-circuit read direction $r_A \propto M^{+} v_A$ (cosine, literal
pseudo-inverse, summed circuit, unit norm). The transpose formula
$r \propto W^{\top} t$, sometimes proposed as the solution of the cosine problem, is
in fact the solution of the dot-product problem; within this crossing it is a
legitimate lever point rather than an approximation.

Finally, the criterion by which the best definition is selected is itself a
definitional choice --- correlational alignment with residual activations at early
tokens, or the causal effect of intervening along $r_A$, and in either case under which
prompt distribution --- and must be fixed before the sweep is run.

\end{document}
